\section{数字格式与区间背景}
\key{value format} 决定显示原值还是计算占比；
\key{number format} 决定数字的排版形式，不改变条长。
\key{fixed} 为默认定点格式，千位使用细空格；\key{grouped} 使用逗号分组；
\key{scientific} 使用科学计数法。\key{precision} 指定小数位数，科学计数法中指尾数的小数位数。

\begin{example}{示例 8：分组数字与科学计数法}
\gradbar[max=200000,
  number format=grouped,
  precision=0]{125000}

\gradbar[max=1,
  number format=scientific,
  precision=2]{0.00001234}
\end{example}

科学计数法也支持 e/E 输入；普通小数输入仍可使用。
极小数不会先按定点精度四舍五入成零；标签较长时应增大 \key{label width}。

\begin{example}{示例 9：参考区间与目标线}
\gradbar[band={60,80},
  band color=lightgreen,
  target=70]{75}

\gradbar[min=-50,max=50,
  band={-10,15},
  band color=skyblue,
  band opacity=0.4]{-25}
\end{example}

\key{band={下限,上限}} 使用与数据相同的单位，要求两个端点严格递增且位于量程内。
背景在轨道上下各延伸 2pt，填充经过区间时仍可看见参考范围。
\key{band color} 默认为 gold，\key{band opacity} 默认为 0.3，取值为 0 到 1。
使用 \key{band={}} 清除区间。

\note{区间背景表示人为给定的参考范围，不代表统计误差。
误差区间仍使用 \key{error}；参考区间的越界端点会直接报错，不自动截断。}
